\subsection{Simplicity of Quaternion Algebras}

PROPOSITION 3. --- Let \(\alpha, \beta, \gamma\) be elements of K, and let \(f\) be a quaternion algebra of type \((\alpha, \beta, \gamma)\) . Denote its Cayley trace and norm by \(\mathrm{T}_{\mathrm{F}}\) and \(\mathrm{N}_{\mathrm{F}}\) . The following properties are equivalent:

\begin{enumerate}
\def\labelenumi{(\roman{enumi})}
\item
  The algebra \(\mathbf{F}\) is central simple.
\item
  For every nonzero element \(x\) of \(\mathrm{F}\) , there exists an element \(y\) in \(\mathrm{F}\) such that \(\mathrm{T}_{\mathrm{F}}(xy) \neq 0\) .
\item
  We have \((4\alpha + \beta^{2})\gamma \neq 0\) .
\end{enumerate}

Suppose that these properties hold. Then for every x in F, the reduced characteristic polynomial of x is \(\mathrm{X}^{2}-\mathrm{T}_{\mathrm{F}}(x)\mathrm{X}+\mathrm{N}_{\mathrm{F}}(x)\) . In particular, \(\mathrm{T}_{\mathrm{F}}(x)\) is the reduced trace of x, and \(\mathrm{N}_{\mathrm{F}}(x)\) is its reduced norm.

(i)⇒(ii): If the algebra F is central simple, then it follows from Proposition 1 of VIII, p.~361 and the definition of the reduced trace (VIII, p.~340, Definition 2) that \(T_{F}\) is its reduced trace; assertion (ii) follows (VIII, p.~343, Proposition 5).

\begin{enumerate}
\def\labelenumi{(\roman{enumi})}
\setcounter{enumi}{1}
\tightlist
\item
  \(\Leftrightarrow\) (iii): Let \((e_i)_{1\leqslant i\leqslant 4}\) be a basis of F of type \((\alpha ,\beta ,\gamma)\) (III, §2, No.~5, p.~445). The matrix \((\mathrm{T_F}(e_i e_j))\) is equal to
\end{enumerate}

\[
\left( \begin{array}{c c c c} 2 & \beta & 0 & 0 \\ \beta & 2 \alpha + \beta^ {2} & 0 & 0 \\ 0 & 0 & 2 \gamma & \beta \gamma \\ 0 & 0 & \beta \gamma & - 2 \alpha \gamma \end{array} \right).
\]

Its determinant is \(-\gamma^{2}(4\alpha + \beta^{2})^{2}\) . The equivalence of properties (ii) and (iii) follows from V, §8, No.~2, p.~49, Lemma 1.

(iii)⇒(i): Suppose \((4\alpha + \beta^{2})\gamma \neq 0\) . We then have \(\gamma \neq 0\) , and we have \(\beta \neq 0\) if K has characteristic 2. By Proposition 2, the algebra F is central. Let x be an element of the Jacobson radical of F. For every \(y \in F\) , xy is nilpotent, so \(T_{F}(xy) = 0\) (Remark 3 of VIII, p.~362). Since (ii) is equivalent to (iii), we have x = 0. This proves that F is a semisimple K-algebra. Since its center is K, it is simple.

The last assertion follows from Proposition 1 of VIII, p.~361 and the definition of the reduced characteristic polynomial (VIII, p.~340, Definition 1).

Denote the characteristic of K by p.~By Proposition 3, if \(p \neq 2\) , then every quaternion algebra over K of type \((\alpha, 0, \gamma)\) with \(\alpha\) and \(\gamma\) in \(K^{*}\) is central simple. If p = 2, then every quaternion algebra of type \((\alpha, 1, \gamma)\) with \(\alpha \in K\) and \(\gamma \in K^{*}\) is central simple. Conversely, we have the following.

PROPOSITION 4. --- Let A be a central simple algebra of degree 4 over K. Denote the characteristic of K by p.

\begin{enumerate}
\def\labelenumi{\alph{enumi})}
\item
  If \(p \neq 2\) , then there exist nonzero elements \(\alpha\) and \(\gamma\) of \(K\) such that the algebra \(A\) is isomorphic to the quaternion algebra of type \((\alpha, 0, \gamma)\) .
\item
  If \(p = 2\) , then there exist an element \(\alpha\) of \(K\) and an element \(\gamma\) of \(K^*\) such that the algebra \(A\) is isomorphic to the quaternion algebra of type \((\alpha, 1, \gamma)\) .
\end{enumerate}

By Wedderburn's theorem (VIII, p.~120, Theorem 1), there exist an integer \(r \geqslant 1\) and a field D with center K such that A is isomorphic to \(\mathbf{M}_r(\mathrm{D})\) . We then have \(r^2[\mathrm{D}:\mathrm{K}] = [\mathrm{A}:\mathrm{K}] = 4\) . If \(r = 2\) , then A is isomorphic to \(\mathbf{M}_2(\mathrm{K})\) , and Proposition 4 follows from the example of VIII, p.~362. Otherwise, we have \(r = 1\) , and A is a field with center K. It then has a maximal commutative subfield E that is a separable extension of K; since

A has degree 4 over K, the extension E has degree 2 over K (VIII, p.~265, Corollary 2). It is therefore quadratic (III, §2, No.~3, p.~439). Let s be the conjugation of E (III, §2, No.~3, p.~440). By the Skolem--Noether theorem (VIII, p.~256, Corollary 1), there exists an invertible element j of A such that we have \(s(x) = jxj^{-1}\) for every x in E. The field E is separable over K, so we have \(s \neq Id_{E}\) , so that \(j \notin E\) . Since A is a vector space of dimension 4 over K, it is a left vector space of dimension 2 over E, so we have \(A = E \oplus Ej\) . We have \(s^{2} = Id_{E}\) , so the element \(j^{2}\) of A belongs to the center of A; hence there exists an element \(\gamma\) of \(K^{*}\) such that \(j^{2} = \gamma\) .

When \(p \neq 2\) , there exist an element i of E and an element \(\alpha \in K^{*}\) such that \(\mathrm{E} = \mathrm{K}(i)\) and \(i^{2} = \alpha\) (V, §11, No.~9, p.~93, Example 3); in this case, A is isomorphic to the quaternion algebra of type \((\alpha, 0, \gamma)\) . When p = 2, there exist an element i of E and an element \(\alpha\) of K such that \(\mathrm{E} = \mathrm{K}(i)\) and \(i^{2} = i + \alpha\) (V, §11, No.~9, p.~93, Example 2), so that A is isomorphic to the quaternion algebra of type \((\alpha, 1, \gamma)\) .

COROLLARY 1. --- Let A be a central simple K-algebra of finite degree \textgreater{} 1 whose elements are all algebraic of degree \(\leqslant 2\) over K. Then A is isomorphic to a quaternion algebra over K.

If K is finite, then the algebra A is isomorphic to a matrix algebra \(\mathbf{M}_{n}(\mathrm{K})\) (VIII, p.~357, Corollary 2) and therefore contains elements of degree n over K; the assumption implies n = 2 and therefore the result in this case (VIII, p.~362, Example). Suppose that the field K is infinite. Let L be a maximal étale subalgebra of A. By V, §7, No.~4, p.~41, Proposition 7, there exists an element x of A such that the K-algebra L is equal to \(K[x]\) , hence by assumption has degree \(\leqslant 2\) . Since we have \([A : K] = [L : K]^{2}\) (VIII, p.~264, Proposition 4 and p.~262, Proposition 3), we conclude that \([A : K] = 4\) . Corollary 1 then follows from Proposition 4.

COROLLARY 2. --- Let \((\mathrm{E}, s)\) be a Cayley algebra over K such that the K-algebra E is central simple of finite degree \textgreater{} 1 over K. Then E is isomorphic to a quaternion algebra over K.

Every element \(u\) of \(\mathrm{E}\) satisfies \(u^2 - \mathrm{T}_{\mathrm{E}}(u)u + \mathrm{N}_{\mathrm{E}}(u) = 0\) , so the K-algebra \(\mathrm{E}\) is isomorphic to a quaternion algebra (Corollary 1).

